% SGA 3, 第 VI_B 篇, 第 1 轮英文重构底本的中译.
% 范围: 命题 7.9 与推论 7.10 全部.
% 权威底本: Polo--Gille Exp6B-13oct24.pdf, 组件第 60 页
% (合卷第 401 页; 重版页眉第 390 页; 原版边栏页码第 374 页).
% 第 8 节始于组件第 60 页稍后位置, 不在本组件内.
% OCR 与 Reinhold 英文 Markdown 仅为对照材料. 本组件无图.

\medskip
\phantomsection\label{prop:sga3-VIB-7.9}%
\noindent\textbf{命题 7.9.---}%
\emph{设 \(k\) 是代数闭域, \(G\) 是局部有限型的 \(k\)-群.}
\begin{enumerate}
\item[\textup{\SGAInlineRefTarget{prop:sga3-VIB-7.9:item:i}(i)}]\phantomsection\label{concl:sga3-VIB-7.9-i}
  \emph{设 \(A\) 与 \(B\) 是 \(G\) 的两个归纳可构造子集.
  以 \(A_0\) 表示属于 \(A\) 的 \(G\) 的有理点所成的集合. 则}
  \[
    (A\cdot B)_0=A_0\cdot B_0,
  \]
  \emph{其中右边的乘积取于群 \(G(k)\) 中.}

\item[\textup{\SGAInlineRefTarget{prop:sga3-VIB-7.9:item:ii}(ii)}]\phantomsection\label{concl:sga3-VIB-7.9-ii}
  \emph{设 \(X\) 是局部有限型的几何约化 \(k\)-概形,
  \(f:X\to G\) 是一个 \(k\)-态射. 置}
  \[
    \Gamma'_G(f)=
      \bigcup_{n\geq1}f^n(X^n).
  \]
  \emph{则 \(\Gamma'_G(f)_0\) 是由 \(f(X)_0\) 生成的 \(G(k)\) 的子群.}

\item[\textup{(iii)}]\phantomsection\label{concl:sga3-VIB-7.9-iii}
  \emph{特别地, 设 \(A\) 与 \(B\) 是 \(G\) 的两个平滑群子概形,
  并按 \hyperref[def:sga3-VIB-7.2.2]{7.2.2} 的记号置}
  \[
    \Gamma'=
      \bigcup_{n\geq1}
        \nu^n\bigl((A\times_k B)^n\bigr)
  \]
  \emph{则 \(\Gamma'_0\) 是 \(A(k)\) 与 \(B(k)\)
  在 \(G(k)\) 中的换位子群.}
\end{enumerate}

\emph{证明.}
先证明 \textup{\SGARefLink{prop:sga3-VIB-7.9:item:i}{(i)}}. 显然有
\[
  A_0\cdot B_0\subset(A\cdot B)_0.
\]
反之, 设 \(z\in(A\cdot B)_0\).
则 \(\mu^{-1}(z)\) 是 \(G\times_kG\) 的一个闭子集,
而 \(A\times_kB\) 是 \(G\times_kG\) 的一个归纳可构造子集
(参见 \hyperref[par:sga3-VIB-3.0]{3.0}). 因此,
\[
  \mu^{-1}(z)\cap(A\times_kB)
\]
是 \(G\times_kG\) 的一个非空归纳可构造子集.
由 EGA~IV\(_3\), 10.4.8, 它因而包含 \(G\times_kG\) 的一个有理点.
其投影 \(x\) 与 \(y\) 是 \(G\) 的有理点,
满足 \(x\in A\), \(y\in B\) 及 \(x\cdot y=z\). 所以
\begingroup
\renewcommand{\thefootnote}{\fnsymbol{footnote}}%
\footnote[1]{\label{note:sga3-VIB-7.9-product}译者注:
法文证明在此印为 \((A,B)_0=A_0\cdot B_0\).
命题陈述及乘法论证要求 \((A\cdot B)_0=A_0\cdot B_0\),
英译稿已据此恢复.}%
\endgroup
\[
  (A\cdot B)_0=A_0\cdot B_0.
\]

要证明 \textup{\SGARefLink{concl:sga3-VIB-7.9-ii}{(ii)}},
注意由于 \(f^n\) 局部有限型, \(f^n(X^n)\)
是 \(G\) 的一个归纳可构造子集
(EGA~IV\(_1\), 1.9.5\textup{(viii)})\begingroup
\renewcommand{\thefootnote}{\fnsymbol{footnote}}%
\footnote[2]{\label{note:sga3-VIB-7.9-EGA-volume}
译者注: 法文原文在此印为 EGA~IV\(_4\).
所引结果 1.9.5\textup{(viii)} 属于 EGA~IV\(_1\),
英译稿已据此恢复.}%
\endgroup.
由断言 \textup{\SGARefLink{prop:sga3-VIB-7.9:item:i}{(i)}} 归纳可知, 若
\[
  A=f(X)_0,
\]
则有
\[
  f^n(X^n)_0=(A\cup A^{-1})^n.
\]
所以,
\phantomsection\label{form:sga3-VIB-7.9-generated-rational-points}%
\[
  \Gamma'_G(f)_0
    =\bigcup_{n\geq1}f^n(X^n)_0
    =\bigcup_{n\geq1}(A\cup A^{-1})^n,
\]
这就是由 \(A=f(X)_0\) 生成的 \(G(k)\) 的子群.

最后, \textup{(iii)} 由 \textup{\SGARefLink{prop:sga3-VIB-7.9:item:ii}{(ii)}}
及定义得到.

\medskip
\phantomsection\label{cor:sga3-VIB-7.10}%
\noindent\textbf{推论 7.10.---}%
\emph{在 \hyperref[prop:sga3-VIB-7.8]{7.8} 的假设下,
若 \(k\) 代数闭, 则 \((A,B)(k)\) 是 \(A(k)\)
与 \(B(k)\) 在 \(G(k)\) 中的换位子群.}

事实上, 只需应用 \hyperref[concl:sga3-VIB-7.9-iii]{7.9\textup{(iii)}},
\hyperref[prop:sga3-VIB-7.8]{7.8} 及 \hyperref[prop:sga3-VIB-7.6]{7.6}.

% 中文翻译检查点：本组件全部正文、证明、脚注及源注已译；源122行。
% 机械核对：行内数学53处、陈列数学9处及全部引用键守恒（数学中的自然语言已译）。
% 状态：完成源文覆盖和逐段自校；尚未经独立合卷编译与页面审查。
